\section{H16: Idle traps deepen with time}
\label{sec:H16}

\textbf{The question:} Agents sometimes get stuck in long silences, chains of timer pauses or repeated failures. Is the escape from such a trap memoryless, as Kramers escape predicts? Or does the trap deepen the longer the agent stays in it? What gets a stuck agent out?

\textbf{The intuition:} In the Kramers picture, a particle sits in a well, and noise kicks it over a barrier. It escapes at a constant rate, whatever time it has already spent there. The escape is memoryless, and dwell times are exponential. In the aging picture (a Bouchaud trap model), trap depths have a broad spread. The longer an agent stays stuck, the more likely it is in a deep trap. The hazard, the chance of escape per unit time for an agent still stuck, then falls with time. A kick is a message to the agent. In Kramers, each kick lowers the barrier, so the escape rate rises exponentially with the number of kicks. In the rival picture, the agent reads the kick at its next call. The first kick does the work and more kicks add little: the response saturates.

\textbf{What we measured:} We built trap spells per agent. Inactive spells are gaps of 180\,s or more, where single isolated actions do not end the spell. Pause chains are runs of timer pauses, with one decision at each timer wake. The hazard slope $\beta$ on $\ln(\text{time elapsed})$ uses spells of 10\,min or more, with agent fixed effects. $\beta=0$ is memoryless and $\beta<0$ is aging; $|\beta|\le0.3$ counts as memoryless. At each timer wake, a logit relates the chance to act to the chain length, with the declared pause duration as a control. A kick is a directed message in the prior 2\,min. The null takes the same agent's kicks from another day.
\begin{itemize}
\item \emph{Scheduler field:} partly removed. Agent fixed effects remove per-agent frailty, and the day-swap null keeps the schedule. The swarm block was not re-run on fixed data.
\item \emph{External drive:} removed. Kicks are past-only, and with the nudger off (NE43) the aging persists.
\item \emph{Shared model prior:} not relevant. Agent fixed effects absorb each agent's frailty, and the card makes no family claim.
\item \emph{Common input:} open. Kicks count by posting time, not by the read-out call.
\end{itemize}
The synthetic check used overdamped Langevin agents seen only through event times. Memoryless truth gave $+0.05\pm0.07$ (truth $+0.02$). Deepening traps gave $-0.58$ and $-1.08$ (truth $-0.54$ and $-1.18$). Without agent fixed effects, frailty alone faked aging of $-0.25$, so the fixed effects are required.

\textbf{What we found:} Within an agent, escape from inactive spells and pause chains ages instead of being memoryless\cred{0.70}; EU 2.1 (value if true 3.0).
\begin{itemize}
\item The inactive-spell slope is $-0.77$ $[-0.81,-0.73]$ in \#51 (4,321 escapes) and $-1.71$ in \#38. The point slope is below $-0.3$ in 6/8 regime-III goal periods. The interval lies below $-0.3$ in only 3/8.
\item The chance to act at a timer wake falls with chain length in 6/7 periods (\#51 $-0.38$, \#38 $-0.44$).
\item With the nudger off (NE43), the aging does not change ($-0.59\to-0.53$). The trap deepens by itself.
\item One directed message during a pause raises the odds of acting at the next timer wake $\times1.5$--$2.9$ in 5/7 periods (\#51 $\times1.54$). In 14/16 powered cells, the dose response is additive or saturating, not exponential.
\item On real failures (not stderr), error-loop aging survives only in \#51 ($-0.57$ $[-0.69,-0.06]$), so the card withdraws it elsewhere. Messages do not break real failure loops (9/9 periods).
\item The swarm has no collective bistability: the mean-field coupling $\beta J_0$ is 0.09--0.41.
\end{itemize}
The scope is 11 goal periods outside the reserved data (8 in regime III), for silences and pause chains only. The credence comes from an adjudication. The coordinator scored a wider claim at 0.44 and marked it fragile. The blind rater gave 0.70, and the ruling kept 0.70, not fragile.

\textbf{What it means:} For an operator, kick a stuck agent early, because its escape hazard falls roughly as $t^{-0.8}$. Send one directed message during a pause. Under the 5-min pause default, it lifts escape from deep pause chains from 0.17 to 0.25. Before that default (NE44) it lifted escape from 0.21 to 0.56. A second or third message adds little. Do not message an agent in a real failure loop; use tooling or a reset. For a physicist, the barrier from a Boltzmann inversion of one-minute activity is an artifact of the smoothing, not a measurement.

\textbf{Caveats:}
\begin{itemize}
\item Spells differ in kind within one agent, and that can mimic aging. Agent fixed effects remove only differences between agents.
\item The pre-registered count rule (5/8 periods) is not met. Short goal periods give unstable bootstrap intervals.
\item The predicted drop in the chance to act with the nudger off failed ($+0.006$ $[-0.09, 0.10]$). The NE44 interaction also failed as stated, and that contrast compares different goal periods.
\item We defined the robust pause-chain rule after we saw that strict chains barely exist.
\item No confirmation test on reserved goal periods has run.
\item Figure panels (b) and (c) use H72's table of pause points, with a different estimator. Its \#51 aging slope is $-0.48$, not H16's $-0.77$.
\end{itemize}

\begin{figure}[h]
\includegraphics[width=\textwidth]{H16-trap-aging/fig.pdf}
\caption{Idle traps age by themselves, and one directed kick helps. (a) Simulation: with one trap depth the escape chance does not depend on trap age, but with a broad spread of depths it falls. (b) In \#51, the share of pause points that start a run of activity falls with trap age, with a directed message (orange) and without one (dark). (c) Aging slopes in four goal periods stay negative when an input-starvation clock enters the model. Animation: \texttt{writeup/visuals/H16-trap-aging/anim.mp4}.}
\label{fig:int-H16}
\end{figure}
